
Our results show that automation and augmentation are not two labels for the same underlying capability. They produce systematically different rankings. The top-ranked model differs between automation and assistant-only augmentation in five of the seven tasks. Thus, the assistant value is sharply task-contingent. Interpreting these patterns requires moving beyond leaderboard thinking to role-specific model selection. The automation results are comparatively concentrated around a small set of
direct solvers, while GPT-3.5-Turbo ranks last. This ordering characterizes performance on our single-turn professional deliverables and should not be interpreted as reproducing the orderings of agentic benchmarks built
around multi-step tool use and environment interaction.
Augmentation tells a different story. No single model dominates across tasks, average ranks are more dispersed, and several strong direct solvers perform poorly as assistants. Most strikingly, on operations research, tax preparation, and travel planning, the unaided GPT-3.5-Turbo worker outranks every assisted condition, and it has the second best average augmentation rank overall after GPT-5-Mini. We read this not as evidence that assistance is useless, but as evidence that its value is conditional and can be negative. Poorly matched or overly complex process guidance can distract the worker, constrain useful execution, or push it away from task requirements.

\subsection{Implications for Organizations and Workflow Design}

Our framework offers several practical applications for individuals and organizations adopting AI systems. First, it can support model selection across teams, products, and workflows. Teams seeking a reliable solver for well-defined, repeatable tasks may prioritize automation performance, whereas teams embedding AI into expert workflows, where the goal is to improve human output rather than replace human judgment, may prioritize augmentation ability.

Importantly, the framework is modular. In our experiments, we use GPT-3.5-Turbo as an illustrative fixed worker model to demonstrate the pipeline and hold the worker role constant across augmentation trials. However, individuals and organizations can substitute this worker model with a model, human baseline, or internal workflow proxy that better reflects their own deployment context. For example, a consulting firm could evaluate how different assistant models guide junior analysts, domain experts, customer-support agents, or a preferred internal model. This flexibility allows the framework to be adapted to local goals, constraints, and performance standards.
    
Secondly, designing human-AI collaboration can be done more efficiently and effectively. Beyond model selection, our framework can guide how organizations structure the division of labor between humans and AI. At the organizational level, different models can be assigned to different subtasks according to where each model's comparative advantage lies. At the team level, human-AI pairs can use the automation/augmentation profiles to decide which subtasks to delegate entirely versus which to treat as collaborative, keeping humans in the loop where an assistance yields the largest gains. At the individual level, a practitioner can use the same logic to choose the tool best matched to their specific workflow, rather than defaulting to a single general-purpose assistant.

\subsection{Implications for Evaluation and Research}

Our findings also speak to how LLMs should be evaluated. Many existing benchmarks measure automation almost exclusively. If gains on those benchmarks do not transfer to augmentation, then labs and leaderboard builders may be optimizing for only one economically relevant capability. Adopting automation and augmentation as distinct evaluation axes would make model profiles more informative for human-in-the-loop and multi-agent deployment. If automation and augmentation are empirically distinct capabilities, then any evaluation that tests only one is incomplete by design, and our framework offers a template for fixing that. Leaderboard builders and AI labs optimizing for benchmark performance could adopt the automation vs. augmentation distinction as a structural requirement for capability evaluation.
    
CentaurBench sits between two literatures. Relative to large CS benchmarks, it evaluates economically grounded occupational tasks under a fixed coaching protocol rather than abstract capability alone. Relative to field experiments, it compares many frontier models across both roles and multiple domains rather than a single tool in a single setting. Relative to marketplace benchmarks such as HAPI \citep{upwork_hapi2025}, it holds the downstream worker fixed to isolate assistant quality across models at scale. The appropriate next step is not to replace one paradigm with another, but to combine them for optimal model selection. 

% rephrased above
%The last application is a meta one. We believe that our framework is itself a critique of how models are currently benchmarked. Researchers building leaderboards, and AI labs deciding what to optimize for, could use the automation/augmentation distinction as a structural requirement for any serious capability evaluation. The implicit argument is that existing benchmarks are incomplete by design, and our framework offers a template for fixing that.
